\documentclass[11pt]{article}
\usepackage[a4paper,margin=25mm]{geometry}
\usepackage[no-math]{fontspec}
\setmainfont{Latin Modern Roman}
\usepackage{amsmath,amssymb,amsthm,xfrac,xspace,hyperref,graphicx,comment}
\excludecomment{DefTralics}
\providecommand{\C}{}
\newcommand{\ie}{i.e.,\xspace}
\newcommand{\eg}{e.g.,\xspace}
\newcommand{\etc}{etc.\xspace}
\newcommand{\cf}{cf.\xspace}
\newcommand{\resp}{resp.\xspace}
\newcommand{\smaller}{\small}
\newcommand{\Subsection}{\subsection}
\newcommand{\Subsubsection}{\subsubsection}
\newcommand{\skpt}{\smallskip}
\emergencystretch=3em
\input{author-preamble.tex}
\begin{document}
\begin{center}{\Large Stable item 1922}\end{center}
\paragraph{Source and attribution.}
Daniel Greb, Matei Toma, \emph{Moduli spaces of sheaves that are semistable with respect to a Kähler polarisation}, Journal de l’École polytechnique — Mathématiques 7 (2020), publisher version, DOI 10.5802/jep.116. \href{https://jep.centre-mersenne.org/articles/10.5802/jep.116/}{Publisher article}. Authors' text: \href{https://creativecommons.org/licenses/by/4.0/}{CC BY 4.0}.
\paragraph{Editorial problem boundary.}
Prove only Theorem 3.2 (Existence) below: for a connected smooth complex projective manifold, a real Kähler polarisation, and torsion-free rank-two sheaves of fixed Chern classes, the stack of Gieseker--Maruyama semistable sheaves admits a good moduli space with affine diagonal. The semistability definition, quotient/slice preliminaries, and the complete two-step existence proof are retained. No assertion of a projective scheme, higher rank, or properness is selected. The remaining article results are conservative author context only.
\paragraph{Difficulty (editorial): H3.}
Rank-two stabiliser classification and equivariant semi-universal slices reduce closed-orbit preservation to a two-weight torus-action contradiction. Extension flag stacks supply the orbit-closure good moduli spaces, allowing the Alper-Fedorchuk-Smyth criterion beyond rational ample polarisations.
\paragraph{Editorial transcription note.}
Original author language, mathematical content and ordering are unchanged. Publisher front matter is replaced by this independent layout wrapper. Original native body and macros remain editable; bibliography is recovered from the matching cached publisher HTML. Adjacent claims are author context, not additional corpus items. Printed mathematical slips are retained without assistant repair.
\clearpage
\input{author-body.tex}
\begin{thebibliography}{99}
\bibitem{AlperFS} J. Alper, M. Fedorchuk \& D. I. Smyth - “Second flip in the Hassett-Keel program: existence of good moduli spaces” , Compositio Math. 153 (2017) no. 8, p. 1584-1609
\bibitem{AHLH} J. Alper, D. Halper-Leistner \& J. Heinloth - “Existence of moduli spaces for algebraic stacks” , 2018
\bibitem{AlperHR} J. Alper, J. Hall \& D. Rydh - “A Luna étale slice theorem for algebraic stacks” , 2015
\bibitem{AlperKresch} J. Alper \& A. Kresch - “Equivariant versal deformations of semistable curves” , Michigan Math. J. 65 (2016) no. 2, p. 227-250
\bibitem{Alper13} J. Alper - “Good moduli spaces for Artin stacks” , Ann. Inst. Fourier (Grenoble) 63 (2013) no. 6, p. 2349-2402
\bibitem{Alper} J. Alper - “Artin algebraization and quotient stacks” , 2015
\bibitem{Ant} S. Antonakoudis - “Criteria of separatedness and properness” , 2019, notes available at https://www.dpmms.cam.ac.uk/\textasciitilde{}sa443/papers/criteria.pdf
\bibitem{BuTeTo1bis} N. Buchdahl, A. Teleman \& M. Toma - “A continuity theorem for families of sheaves on complex surfaces” , J. Topology 10 (2017) no. 4, p. 995-1028
\bibitem{D_Luna} J.-M. Drézet - “Luna’s slice theorem and applications” , in Algebraic group actions and quotients , Hindawi Publ. Corp., Cairo, 2004, p. 39-89
\bibitem{Eisenbud_CommAlg} D. Eisenbud - Commutative algebra. With a view toward algebraic geometry , Graduate Texts in Math. , vol. 150 , Springer-Verlag, New York, 1995
\bibitem{FlennerHabil} H. Flenner - Deformationen holomorpher Abbildungen , Osnabrücker Schriften zur Mathematik (Reihe P) , vol. 8 , Fachbereich Mathematik, Univ. Osnabrük, 1978, available online at http://www.ruhr-uni-bochum.de/imperia/md/content/mathematik/lehrstuhli/deformationen.pdf
\bibitem{Gieseker} D. Gieseker - “On the moduli of vector bundles on an algebraic surface” , Ann. of Math. (2) 106 (1977) no. 1, p. 45-60
\bibitem{GRTsurvey} D. Greb, J. Ross \& M. Toma - “Moduli of vector bundles on higher-dimensional base manifolds—construction and variation” , Internat. J. Math. 27 (2016) no. 7, article ID 1650054, 27 pages
\bibitem{GRT1} D. Greb, J. Ross \& M. Toma - “Variation of Gieseker moduli spaces via quiver GIT” , Geom. Topol. 20 (2016) no. 3, p. 1539-1610
\bibitem{GrebToma} D. Greb \& M. Toma - “Compact moduli spaces for slope-semistable sheaves” , Algebraic Geom. 4 (2017) no. 1, p. 40-78
\bibitem{Hartshorne} R. Hartshorne - Algebraic geometry , Graduate Texts in Math. , vol. 52 , Springer-Verlag, New York-Heidelberg, 1977
\bibitem{Hartshorne_Deformations} R. Hartshorne - Deformation theory , Graduate Texts in Math. , vol. 257 , Springer, New York, 2010
\bibitem{HL} D. Huybrechts \& M. Lehn - The geometry of moduli spaces of sheaves , Cambridge Math. Library , Cambridge University Press, Cambridge, 2010
\bibitem{HMP} P. Heinzner, L. Migliorini \& M. Polito - “Semistable quotients” , Ann. Scuola Norm. Sup. Pisa Cl. Sci. (4) 26 (1998) no. 2, p. 233-248
\bibitem{JoyceSong} D. Joyce \& Y. Song - A theory of generalized Donaldson-Thomas invariants , Mem. Amer. Math. Soc. , vol. 217, no. 1020 , American Mathematical Society, Providence, RI, 2012
\bibitem{Kempf} G. R. Kempf - Algebraic varieties , London Math. Society Lect. Note Ser. , vol. 172 , Cambridge University Press, Cambridge, 1993
\bibitem{Knutson} D. Knutson - Algebraic spaces , Lect. Notes in Math. , vol. 203 , Springer-Verlag, Berlin-New York, 1971
\bibitem{KosarewStieber90} S. Kosarew \& H. Stieber - “A construction of maximal modular subspaces in local deformation theory” , Abh. Math. Sem. Univ. Hamburg 60 (1990), p. 17-36
\bibitem{Langton} S. G. Langton - “Valuative criteria for families of vector bundles on algebraic varieties” , Ann. of Math. (2) 101 (1975), p. 88-110
\bibitem{Lange} H. Lange - “Universal families of extensions” , J. Algebra 83 (1983) no. 1, p. 101-112
\bibitem{LaumonMB} G. Laumon \& L. Moret-Bailly - Champs algébriques , Ergeb. Math. Grenzgeb. (3) , vol. 39 , Springer-Verlag, Berlin, 2000
\bibitem{LePotier_Lectures} J. Le Potier - Lectures on vector bundles , Cambridge Studies in Advanced Math. , vol. 54 , Cambridge University Press, Cambridge, 1997
\bibitem{Palamodov-versal} V. P. Palamodov - “Deformations of complex spaces” , in Several complex variables. IV. Algebraic aspects of complex analysis (G. M. Khenkin, ed.) , Encycl. Math. Sci. , vol. 10 , Springer-Verlag, Berlin, 1990, p. 105-194
\bibitem{Seshadri} C. S. Seshadri - “Space of unitary vector bundles on a compact Riemann surface” , Ann. of Math. (2) 85 (1967), p. 303-336
\bibitem{stacks-project} Stacks Project Authors - “The Stacks Project” , 2019, https://stacks.math.columbia.edu
\bibitem{TelemanCommContMath} A. Teleman - “Families of holomorphic bundles” , Commun. Contemp. Math. 10 (2008) no. 4, p. 523-551
\bibitem{TomaLimitareaI} M. Toma - “Bounded sets of sheaves on Kähler manifolds” , J. reine angew. Math. 710 (2016), p. 77-93
\bibitem{TomaCriteria} M. Toma - “Properness criteria for families of coherent analytic sheaves” , 2017, to appear in Algebraic Geom.
\bibitem{TomaLimitareaII} M. Toma - “Bounded sets of sheaves on Kähler manifolds. II” , 2019
\end{thebibliography}
\end{document}
